\clearpage
\fancyhead[L]{\small COMET.}
\SourceLeaf{447}
\begin{center}
{\large COMET.}

\noindent\rule{\linewidth}{0.4pt}

Göttingische gelehrte Anzeigen. No. 26. pp. 249..251. 15 February 1823.

\noindent\rule{\linewidth}{0.4pt}
\end{center}

A letter from Mr. \textsc{Rümker} at Paramatta in New South Wales to Councillor \textsc{Gauss}, dated August 1822, contains the observations of \textsc{Encke}’s comet that Mr. \textsc{Rümker} made in June last year, along with many other astronomical observations made there that are being held for publication elsewhere. As is well known, there was very little hope that the comet could be observed in Europe on its return to the Sun in 1822, and in fact all the efforts of European astronomers to see it proved entirely fruitless. On the other hand, the expectation was that it would not escape observers in the southern hemisphere. But as word had already come that the English astronomers’ searches at the Cape of Good Hope had \emph{not succeeded}, news from New Holland was awaited all the more eagerly. We are therefore all the more pleased to report in our journal the confirmation obtained by our able compatriot. It was in no. 83 of 1819 [p.~420 of this volume.] that we first reported Mr. \textsc{Encke}’s most remarkable discovery. The following are \emph{all} the observations made by Mr. \textsc{Rümker}. After 23 June the moonlight interfered, and after the full Moon the comet was too faint for any further observations.

\Needspace{23\baselineskip}
\begin{center}
Observations of \textsc{Encke}’s comet at \emph{Paramatta}.
\end{center}

\begin{center}
\small\setlength{\tabcolsep}{3pt}
\begin{tabular}{l|ccc|ccc|ccc}
\multicolumn{1}{c|}{1822} & \multicolumn{3}{c|}{Sidereal time} & \multicolumn{3}{c|}{Mean R.A.} & \multicolumn{3}{c}{Mean declination} \\
\hline
June 2 & $10^{\mathrm{u}}$ & $39^{\mathrm{m}}$ & $25^{\mathrm{s}}$ & $92^{\circ}$ & $43'$ & $51''3$ & $17^{\circ}$ & $39'$ & $46''3$ N. \\
3 & 11 & --- & --- & 93 & 46 & 20.7 & 16 & 53 & 7.5 \\
4 & 11 & 3 & 0 & 94 & 46 & 0.0 & 16 & 4 & 36.7 \\
6 & 11 & 7 & 38 & 96 & 42 & 11.6 & 14 & 22 & 42.0 \\
7 & 11 & 3 & 10 & 97 & 38 & 15 & 13 & 26 & 5 \\
8 & 11 & 17 & 25 & 98 & 33 & 47.7 & 12 & 31 & 18.6 \\
10 & 11 & 20 & 0 & 100 & 24 & 43.8 & 10 & 29 & 49.5 \\
11 & 11 & 24 & 39 & 101 & 19 & 44.5 & 9 & 26 & 4.6 \\
12 & 11 & 40 & 0 & 102 & 17 & 52 & 8 & 18 & 30 \\
13 & 11 & 42 & 4 & 103 & 15 & 2 & 7 & 6 & 30 \\
14 & 11 & 55 & 0 & 104 & 15 & 40 & 5 & 52 & 27 \\
15 & 11 & 40 & 48 & 105 & 17 & 0.5 & 4 & 33 & 40 \\
19 & 12 & 13 & 38 & 109 & 54 & 36.4 & 1 & 29 & 43.7 S. \\
20 & 12 & 16 & 53 & 111 & 14 & 26.9 & 3 & 14 & 29.1 \\
22 & 13 & 18 & 46 & 114 & 12 & 20.5 & 7 & 8 & --- \\
23 & 12 & 53 & 55 & 115 & 47 & 41.7 & 9 & 9 & 48.4 \\
\end{tabular}
\end{center}

The report does not say whether \emph{mean}, as applied to right ascension and declination, refers to an average of several observations or to the removal of nutation. Stars from \textsc{Piazzi}’s catalogue and from the \emph{Histoire Céleste} were used for comparison. Mr. \textsc{Rümker} promises to provide further details at the first opportunity.

From the two solstices of December 1821 and June 1822, Mr. \textsc{Rümker} finds the altitude of the pole at the Paramatta observatory to be

\[
33^{\circ}\,48'\,41''97
\]

\Needspace{16\baselineskip}
From an occultation of $\alpha$ Scorpii on 10 April 1822, when

% 
\SourceLeaf{448}
\begin{center}
\small
\begin{tabular}{l l}
ingress \mbox{.\,.\,.\,.\,.\,.\,.\,.\,.\,.\,.} & $18^{\mathrm{u}}\;35^{\mathrm{m}}\;47^{\mathrm{s}}4$ Mean time \\
egress \mbox{.\,.\,.\,.\,.\,.\,.\,.\,.\,.\,.\,.} & $19\quad14\quad27.9$
\end{tabular}
\end{center}

\noindent were observed, Mr. \textsc{Rümker}, using the Moon’s positions from the Nautical Almanac, finds the longitude from Paris to be
\begin{center}
\begin{tabular}{l l}
10 St. $3^{\mathrm{m}}\;56^{\mathrm{s}}3$ & from the ingress \\
$10\quad4\quad7.3$ & from the egress.
\end{tabular}
\end{center}
